% ID: SEL-D0006-OB-E02
% ESTADO: cotejada. Segunda pasada visual completa y borrador compilado revisados.
% Fuente: Selectividad/sel_2019_matematicas_aplicadas/Reserva B/Exámenes Matemáticas Aplicadas CCSS - A y B.pdf
% SHA-256: f9236ec66af0dbd3fbd0231fdcaa6d922425c04d48178ea0737073e76e8eb7d5
% Página: 2; opción B; ejercicio 2. Todos los apartados, incluida la integral.

Se considera la función $f(x)=\dfrac{1}{3}x^3-2x^2+3x+1$.
\begin{enumerate}[label=\alph*)]
  \item \textbf{(0.75 puntos)} Estudie su monotonía y halle sus extremos
  relativos.
  \item \textbf{(0.75 puntos)} Determine los intervalos de concavidad y
  convexidad. Calcule su punto de inflexión.
  \item \textbf{(0.5 puntos)} Calcule la pendiente de la recta tangente a
  la gráfica de $f$ en el punto de abscisa $x=0$.
  \item \textbf{(0.5 puntos)} Calcule $\displaystyle\int f(x)\,dx$.
\end{enumerate}
